\documentclass{article}
\usepackage{amsmath,amssymb}
\title{A non-small boundary interaction left by the complete weak-coupling cube}\author{}\date{}
\begin{document}\maketitle
Use exactly the complete conditional Wilson block of the complete conditional cube calculation. Its normalized fluctuation law has a small error relative to a pinned Gaussian, but its induced boundary action is not a weak interaction relative to independent boundary links. The following coefficient makes this distinction explicit.

Choose a free cube edge $e=(x,\mu)$, with $\mu\in\{1,2,3\}$. Let $e_+$ and $e_-$ be its parallel translates by one unit in the positive and negative fourth directions. Each is a bath link occurring in precisely one of the 54 plaquettes in the complete conditional record: the vertical plaquette joining it to $e$. No original plaquette contains both $e_+$ and $e_-$.

Vary only those two bath links in one fixed Lie direction $T$, normalized by $|T|=1$. Choose the signs of parameters $t_+,t_-$ so their contributions to the corresponding linearized plaquette curls are $A_e+t_+$ and $A_e+t_-$. Every other bath link remains the identity. The exact real Wilson action has quadratic part
\[
 S_b(A)=\frac1{2N}\left[
 A^TH_0A+2A_e(t_++t_-)+t_+^2+t_-^2
 \right]+O((|A|+|t|)^3),\quad H_0=4I+d^Td.
 \tag{1}
\]
All other Lie components decouple at this order. No Gaussian approximation is used to identify this exact Taylor coefficient of the nonlinear action.

The cube symmetry makes every diagonal entry of $H_0^{-1}$ equal. Its verified spectrum from the complete conditional cube calculation therefore gives
\[
 (H_0^{-1})_{ee}=\frac1{12}\left(\frac74+\frac38+\frac2{10}\right)=\frac{31}{160}.
 \tag{2}
\]
Minimizing the quadratic part, or differentiating the exact nearby minimizer by the implicit function theorem, yields
\[
 S_b(U_*(b))=\frac1{2N}\left[
 t_+^2+t_-^2-\frac{31}{160}(t_++t_-)^2
 \right]+O(|t|^3).
 \tag{3}
\]
Thus the exact classical block action has the mixed boundary Hessian
\[
 \partial_{t_+}\partial_{t_-}S_b(U_*(b))\big|_{t=0}
 =-\frac{31}{160N}.
 \tag{4}
\]

For completeness, the finite-dimensional Laplace estimate in the complete conditional cube calculation holds in every fixed $C^r$ norm on a smaller closed bath neighborhood. After shifting to the smooth minimizer, Taylor remainder estimates and their bath derivatives remain uniform; differentiating the rescaled Gaussian integral produces fixed polynomial factors, which are integrable under the same strict Gaussian bound. Differentiated exterior integrals retain the compact action gap. This proves a $C^r$ error of order $\beta^{-1/2}$ for fixed $r$ without any number-of-blocks or all-history assertion.

Taking two bath derivatives of its scalar formula, one obtains the actual induced action coefficient
\[
 \partial_{t_+}\partial_{t_-}[-\log D_\beta(b)]\big|_{t=0}
 =-\frac{31\beta}{160N}
 +\frac12\partial_{t_+}\partial_{t_-}\log\det H_b\big|_{t=0}
 +O(\beta^{-1/2}).
 \tag{5}
\]
The determinant derivative is a finite $\beta$-independent quantity. The original bath-only Wilson action has zero mixed derivative in these two links, since no plaquette contains both. Hence (5) is a genuine induced interaction in the exact partially integrated lattice action, not an artifact of dropping an existing cross term.

On fluctuation-size boundary variables $t_\pm=y_\pm/\sqrt\beta$, the leading mixed coefficient is the nonzero constant $-31/(160N)$. Thus it is not $O(\beta^{-1/2})$ even in that natural scaled boundary chart. It must be retained in the new quadratic/classical action. The determinant term and $6d_G\log\beta$ scalar from \texttt{complete\_conditional\_cube.tex} likewise remain part of the complete carrier; the logarithmic scalar is independent of $t$ but not bounded as $\beta$ increases.

This calculation indicates the viable constructive organization: retain the Schur-complement Gaussian and the full classical minimum as the next reference, then control the nonlinear residual. A small residual on one pinned block does not prove uniform invertibility or summability of those references after many blocks, especially in the global harmonic/gauge directions identified in \texttt{complete\_conditional\_cube.tex}. No estimate of complete S94 follows without that further control and the original selector/record identification.
\end{document}
